A \emph{walk} in a digraph is a finite sequence $(v_0, v_1, \ldots,
v_r)$, $r\geq 1$, of vertices such that $(v_i, v_{i + 1})$ is an arc for all
$i\in \{0, \ldots, r - 1\}$; the \emph{length} of this walk is $r$. A
\emph{path} is a walk where all vertices are distinct. A \emph{cycle} in a
digraph is a walk where $v_0 = v_r$ and all other vertices are distinct.
A digraph is called \emph{acyclic} if it has no cycles.

A \emph{graph} $G$ is defined to be a digraph where $(u, v)$ is an arc if and
only if $(v, u)$ is an arc in $G$. We refer to the pair of arcs above as the
\emph{edge} $\{u, v\}$. Vertices $u$ and $v$ of a graph $G$ are
\emph{adjacent} if $\{u, v\}$ is an edge of $G$. 

An induced subdigraph of a graph, is also a graph, which we refer to as an
\emph{induced subgraph}. A \emph{spanning subgraph} (as opposed to an
induced subgraph) of a graph $G = (V, A)$ is any graph $H = (V, B)$ where $B
\subseteq A$. 

The \emph{degree} of a vertex in a graph is its in-degree, which equals
its out-degree. If $u$ and $v$ are vertices of a graph $G$, then the
\emph{distance} from $u$ to $v$ is the length of a shortest path from $u$ to
$v$, if such a path exists.